% !TEX TS-program = xelatex
% !TEX encoding = UTF-8 Unicode
\documentclass[11pt]{article}
\usepackage[margin=1in]{geometry}

\usepackage{xcolor}
\usepackage{hyperref}
\hypersetup{
  colorlinks=true,
  citecolor=violet
}

\usepackage[style=apa]{biblatex}
\addbibresource{references.bib}

\usepackage{fontspec} 
\setmainfont{Charis SIL}


\title{Thesis Pre-proposal: Matt}
\author{Wren Arbor}
\date{August 15th, 2025}

\begin{document}
\maketitle

If I were to be assigned Matt as my thesis advisor, I would investigate direct/inverse alignment systems, with primary focus on the direct/inverse aspects of Ichishkíin (Sahaptin), a native language of the pacific northwest. Research questions might include:  What is the relationship of direct/inverse to case and ergativity? Is deixis or topicality more important to direct/inverse? What can account for the non-prototypical nature of Ichishkíin's complex direct/inverse system? This would build off of my final project in Neda's Pacific Northwest Languages Topics in Spring 2023, in which I reviewed the direct/inverse system of Ichishkíin and compared it against the prototypical direct/inverse system of Plains Cree in order to form an understanding of how both languages can be seen to be exhibiting hierarchical direct/inverse behavior, and provided a (limited) analysis of the standing of various person-marking morphemes within the case and direct-inverse system. In that project, the data on Ichishkíin was primarily from the grammar by \textcite{jansen2010}; for data on Plains Cree, it used \textcite{dahlstrom1986}. There, I ran into confusion between direct-inverse systems to nominal-hierarchy-based split-ergative system (a similarity discussed in \cite{delancey_interpretation_1981}). Studying this system in closer detail to analyse whether Ichishkíin is actually better described as a split-ergative language would be illuminating in explaining its differences from a more prototypically direct-inverse language like Plains Cree, where inverseness is uniformly marked on the verb with the same marker across situations. I also found that direct-inverse functions on two levels, one of which is primarily based on deixis and the distinction of speech act participants from third persons (as considered primary by \cite{zúñiga2006}), and one of which is more concerned with the distinction of different types of third persons based on pragmatic context (including topicality) or other hierarchies such as animacy (considered primary by \cite{givón1994}). Although it is not necessarily the case that one or the other of these subsystems must be primary, a unified analysis based on a single tendency would be preferable; collecting data on which subsystem is more important to Ichishkíin would help to answer this question.

Methodologically, as far as data collection, collecting new data from existing native speakers would be ideal, and may be possible as there are speakers in Oregon. In order to do this, I would likely want to first make connection with Joana Jansen, who worked on Ichishkíin for her PhD and now is the Associate Director of the Northwest Indian Language Institute at the Univeristy of Oregon. I would ask for her help and guidance with entering the community, to which she likely remains in connection with, and conducting fieldwork. Analytically, I would then use skills learned when I took Field Methods with Sameer to organize, transcribe, and analyse my data, in order to form an analysis with a wealth of recorded examples. If making contact with the community and collecting data directly is would not turn out to be feasible, it is also possible I could use existing texts, as well as existing data in the literature, to do my work.

This proposal, although of interest to me, is also my least preferred, as my knowledge and grasp of morphosyntax is limited, having taken MDA instead of MST; I would likely struggle to engage with the theory quite as deeply as I could in my sociolinguistics proposal for Allison.

\printbibliography
\end{document}